\chapter{Tuổi Học Trò Và Các Cú Ngố Nghê}
\chapterintro{Khịa học sinh mà không giữ lại một chương cho những pha ngố chung của cả tuổi học trò thì chưa trọn}{Đây là phần cười mềm, nhìn lại để thương nhiều hơn chê.}

\begin{lucbatbox}{Lục bát: Lớn mà còn vụng}
Tuổi này tưởng đã đủ khôn\
Mà còn lắm lúc dại hơn trẻ con\
Lúc thì quên vở quên bài\
Lúc thì nhớ nhầm những điều vu vơ\
Nhưng mà nhờ những ngố khờ\
Mai sau nhìn lại mới vừa thấy thương
\end{lucbatbox}

\begin{tutuyetbox}{Thất ngôn tứ tuyệt: Ngố mà đáng nhớ}
Học trò đôi lúc ngố ghê\
Làm sai một chút đã về bâng khuâng\
Nhưng nhờ mấy cú lưng chừng\
Ký ức sau đó sáng bừng hơn lên
\end{tutuyetbox}

\begin{ghichu}
Khịa cho vui, nhưng đến chương cuối thì nên để lại một chút dịu dàng: tuổi học trò có thể dở, nhưng chính cái dở đó làm nó đáng nhớ.
\end{ghichu}

\ngancach

\begin{lucbatbox}{Lục bát: Hôm nay lại ngố}
Hôm nay em lỡ một điều\
Ngày mai lại ngố thêm nhiều chút thôi\
Bạn bè cười rộ một hồi\
Rồi đâu cũng hóa thành lời kể vui\
Nếu mà tuổi trẻ cứ xuôi\
Không vài cú ngố, chắc rồi thiếu duyên
\end{lucbatbox}

\begin{tutuyetbox}{Thất ngôn tứ tuyệt: Cười mình trước đã}
Mai lớn nhìn lại ngày xưa\
Sẽ thương cái vụng còn thừa nơi ta\
Nếu hôm nay biết cười qua\
Con đường đi tiếp cũng hòa hơn thêm
\end{tutuyetbox}